% Template for PLoS
% Version 3.4 January 2017
\documentclass[10pt,letterpaper]{article}
\usepackage[top=0.85in,left=2.75in,footskip=0.75in]{geometry}

% amsmath and amssymb packages, useful for mathematical formulas and symbols
\usepackage{amsmath,amssymb}

% Use adjustwidth environment to exceed column width (see example table in text)
\usepackage{changepage}

% Use Unicode characters when possible
\usepackage[utf8x]{inputenc}

% textcomp package and marvosym package for additional characters
\usepackage{textcomp,marvosym}

% cite package, to clean up citations in the main text. Do not remove.
% \usepackage{cite}

% Use nameref to cite supporting information files (see Supporting Information section for more info)
\usepackage{nameref,hyperref}

% line numbers
\usepackage[right]{lineno}

% ligatures disabled
\usepackage{microtype}
\DisableLigatures[f]{encoding = *, family = * }

% color can be used to apply background shading to table cells only
\usepackage[table]{xcolor}

% array package and thick rules for tables
\usepackage{array}

% create "+" rule type for thick vertical lines
\newcolumntype{+}{!{\vrule width 2pt}}

% create \thickcline for thick horizontal lines of variable length
\newlength\savedwidth
\newcommand\thickcline[1]{%
  \noalign{\global\savedwidth\arrayrulewidth\global\arrayrulewidth 2pt}%
  \cline{#1}%
  \noalign{\vskip\arrayrulewidth}%
  \noalign{\global\arrayrulewidth\savedwidth}%
}

% \thickhline command for thick horizontal lines that span the table
\newcommand\thickhline{\noalign{\global\savedwidth\arrayrulewidth\global\arrayrulewidth 2pt}%
\hline
\noalign{\global\arrayrulewidth\savedwidth}}


% Remove comment for double spacing
%\usepackage{setspace}
%\doublespacing

% Text layout
\raggedright
\setlength{\parindent}{0.5cm}
\textwidth 5.25in
\textheight 8.75in

% Bold the 'Figure #' in the caption and separate it from the title/caption with a period
% Captions will be left justified
\usepackage[aboveskip=1pt,labelfont=bf,labelsep=period,justification=raggedright,singlelinecheck=off]{caption}
\renewcommand{\figurename}{Fig}

% Use the PLoS provided BiBTeX style
% \bibliographystyle{plos2015}

% Remove brackets from numbering in List of References
\makeatletter
\renewcommand{\@biblabel}[1]{\quad#1.}
\makeatother

% Leave date blank
\date{}

% Header and Footer with logo
\usepackage{lastpage,fancyhdr,graphicx}
\usepackage{epstopdf}
\pagestyle{myheadings}
\pagestyle{fancy}
\fancyhf{}
\setlength{\headheight}{27.023pt}
\lhead{\includegraphics[width=2.0in]{PLOS-submission.eps}}
\rfoot{\thepage/\pageref{LastPage}}
\renewcommand{\footrule}{\hrule height 2pt \vspace{2mm}}
\fancyheadoffset[L]{2.25in}
\fancyfootoffset[L]{2.25in}
\lfoot{\sf PLOS}

%% Include all macros below
\newcommand{\lorem}{{\bf LOREM}}
\newcommand{\ipsum}{{\bf IPSUM}}





\usepackage{forarray}
\usepackage{xstring}
\newcommand{\getIndex}[2]{
  \ForEach{,}{\IfEq{#1}{\thislevelitem}{\number\thislevelcount\ExitForEach}{}}{#2}
}

\setcounter{secnumdepth}{0}

\newcommand{\getAff}[1]{
  \getIndex{#1}{York University}
}

\providecommand{\tightlist}{%
  \setlength{\itemsep}{0pt}\setlength{\parskip}{0pt}}

\begin{document}
\vspace*{0.2in}

% Title must be 250 characters or less.
\begin{flushleft}
{\Large
\textbf\newline{Simple principles for engineering reproducible solutions to
environmental management challenges.} % Please use "sentence case" for title and headings (capitalize only the first word in a title (or heading), the first word in a subtitle (or subheading), and any proper nouns).
}
\newline
\\
Christopher J. Lortie\textsuperscript{\getAff{Biology, York University}}\textsuperscript{*},
Malory Owen\textsuperscript{\getAff{York University}}\\
\bigskip
\textbf{\getAff{York University}}Biology, 4700 Keele St.~Toronto, ON, Canada, M3J1P3\\
\bigskip
* Corresponding author: lortie@yorku.ca\\
\end{flushleft}
% Please keep the abstract below 300 words
\section*{Abstract}
An environmental management challenge is an opportunity to use
fundamental science to inform evidence-based decisions for environmental
stakeholders and conservationists. Contemporary science is embracing
open science and increasingly conscious of reproduciblility.
Synergistically, applying these two paradigms in concert advances our
capacity to move beyond context dependency and singlular thinking to
reverse engineer solutions from published scientific evidence associated
with one challenge to many. Herein, we provide a short list of
principles that can guide those that seek solutions to address
environmental management through primary scientific literature.

% Please keep the Author Summary between 150 and 200 words
% Use first person. PLOS ONE authors please skip this step.
% Author Summary not valid for PLOS ONE submissions.
\section*{Author summary}
Grand challenges require grand solutions. Environmental management
cannot neglect fundamental science as a substrate for effective decision
making, and scientists should be conscious of how their science can be
used by managers.

\linenumbers

% Use "Eq" instead of "Equation" for equation citations.
\section{Introduction}\label{introduction}

Conservation decisions can reside with legislators (indirect impactors)
or with environmental managers (direct impactors). To focus on the
latter, managers typically have scientific backgrounds and routinely
navigate the technical literature. However, engagement with scientific
literature is non-trivial for all practitioners because of time, access,
relevance of the science, and reporting standards. Environmental
managers and conservationists need to be able to use evidence to inform
decisions {[}1,2{]}. However, there can be a gap in communication
between basic science and management for at least three reasons.
Firstly, the research is not a direct study of an ecosystem, and an
immediate, real-world solution is needed by managers - preferably with a
demonstrable outcome and reasonable cost estimate {[}3,4{]}. This is a
very real limitation in the primary science literature restoration
ecology for instance {[}5{]}. Secondly, the link between the biology or
ecology present in the literature is not articulately connected to the
similar process for the system at hand. There are notable examples with
journals just as the Journal of Applied Ecology, Basic and Applied
Ecology, the Journal of Environmental Engineering, and many others.
Nonetheless, there remains a vast pool of opportunity for solution
development from publications in other journals. Context-specific
findings in science are a legitimate and useful means to advance
discovery, but at times, studies from one system can be re-purposed for
insights into another. Finally, the capacity to see the forest for the
trees or the generality of a specific, even large-scale or broad basic
research study can be a challenge. Science can be very specialized
{[}6{]}, and mobilizing knowledge for solutions requires both detailed
expertise, scientific synthesis tools {[}7{]}, or a focus on identifying
the salient elements associated with a study {[}8,9{]}. Often, seeing
the forest also requires sampling many trees. This leads to the general
proposal here that experts can promote reproducible and mobile knowledge
for many fundamental science endeavors by considering these principles.
This is both a set of principles for how to make your research
potentially reusable by environmental managers and conservationists
{[}10{]} and inform solutions for the environmental crisis.

An environmental management challenge is a problem presented in
scientific literature or society that, when redefined and reviewed using
these 10 principles, can result in a solution to the original problem.
EXPAND\ldots{} a tiny bit more..how. Typically, a challenge is\ldots{}.
look up how they define grand challenges and make mention here including
the 17 global sustainability goals. A (reverse-engineered) reproducible
solution is a suggested solution to a challenge derived from identifying
all the components of the challenge. It should be applicable to multiple
local-extent challenges when tweaked to fit the circumstances. Finally,
a tool or solution is the desired outcome from the primary research to
support evidence-based decision making in conservation. In this case, a
tool is a methodology researchers use that can facilitate managers to
either identify best ways to measure/identify issues or to provide
solutions for their specific challenge. Any tool is linked to its
respective reproducible solution by the fundamental concepts of
reproducibility {[}11{]}. These can include primarily conceptual
replication, i.e.~repeating the ideas, but there are many other
solutions. Here, we propose that both direct replication (replicating
the same approach in another context) and conceptual replication
(repeated tests of the same concept but with different methods) {[}12{]}
will advance our capacity to explore reproducibility of basic science to
different challenges associated with environmental management. The
primary goal is to escape the everything is context-specific assumption
in many natural science subdisciplines.

The heuristic developed here was inspired by the ``Ten Simple Rules''
format pioneered by Phillip Bourne in the field of computational biology
{[}13{]}. We propose that by distilling the concepts that promote
engagement with scientific literature outside of the research community,
managers can rely on broader sources of scientific knowledge to make
decisions. Additionally, researchers can better understand the
perspective of managers facilitating science and scientific
communication that is more applicable to managers without compromising
their research programs. Here, we will outline and discuss simple
``principles'' scientists can use to make their research more applicable
to managers, and managers can use to identify basic science that fits
their needs.

\section{The principles}\label{the-principles}

\textbf{1. Reframe the problem as challenge.} Doom-and-gloom is a
pervasive mentality in ecology and environmental sciences which can
reduce productivity and problem solving. It can shut down even the most
motivated of minds--but beyond the issue of motivation, reframing a
problem as a \emph{challenge} can reveal solution sets that may
otherwise remain hidden. For example, consider the problem of
human-wildlife conflict between carnivores and the people living near
the Ruaha National Park boundary in Tanzania. The \emph{problem} is that
98.5\% of people perceive wildlife as a threat to their livestock
resulting in increased likelihoods for human-wildlife conflict {[}14{]}.
Re-framed, the challenge can be to decrease the negative perception of
wildlife in areas with high human-wildlife conflict opportunities. It is
a small change in semantics, but it is a potentially profound change in
direction.

\textbf{2. Describe the scope and extent of the challenge.} Defining the
scope of a challenge conceptually and the extent geographically will
ensure that potential solutions fit the challenge. Moving across scales
is a common issue in ecology {[}15{]}, and proposing a spatial scale,
using common terms, and describing the breadth of the challenge will
accelerate interdisciplinary solutions (i.e.~the wildlife-human
challenge above is ecological and societal). The challenge can be
problematic on multiple scales including local, regional, or global
scale and solutions can be needed for each. Conceptually, the scope is
broad in the human-wildlife conflict example whilst the extent is
primarily local to the area surrounding the Southern border of the Ruaha
National Park.

\textbf{3. Explicitly link the basic science to management implications
and policy.} Perhaps the most facile principle, a simple description and
definition of the basic scientific evidence in a study and how it can be
reused is a fundamental step in linking science to evidence-based
decision making for environmental challenges. In the wildlife-human
challenge, depredation of livestock impacted 61.1\% of households, but
livestock loss due to disease or theft was X percent? of total loss
{[}14{]}. Consequently, a clear statement of evidence can illuminate the
most viable solution sets in some instances.

\textbf{4. Propose implications of ignoring this challenge.} A
description of the impact a challenge on a system if left unchecked will
help clarify the severity of the challenge. The trickle down effects and
indirect implications of the challenge should also be examined. For
instance, anti-carnivore sentiment will likely only grow as climate
change and pressures to confine pastoral herders makes livestock more
difficult to raise {[}16,17{]}. Many large carnivores are already
threatened/endangered, and further anthropogenic pressures on the
populations will lead to severe declines in populations including
potential extinction of keystone species {[}18--20{]}; but it is often
associated with underlying human-human conflict {[}21{]}.

\textbf{5. State the direct human needs associated with this challenge.}
State the direct needs of humans as part of the process of generating
reproducible solutions for environmental challenges. The intrinsic value
of the ecosystem is impossible to quantify {[}22{]}, but linking the
challenge and its solutions to direct human needs makes it less likely
to be dismissed. Identifying anthropogenic needs will help a problem
solver create a solution that is appropriate for the challenge, and it
can also prevent the emergence of new related challenges or pressures on
the system in question. This principle can also include engagement with
stakeholders as a mechanism to inform benefits and solutions
{[}23,24{]}.

\textbf{6. List at least one limitation of the study and explain.} There
is no perfect experiment {[}25{]} or synthesis {[}26{]}. Critically
reading the study associated with the challenge can mean the difference
between success and failure of a later implemented management solution
that otherwise follows all other principles presented here. A clear
analysis of causation and correlation can help avoid a fatal misstep and
ensures effective framing of expected outcomes with an environmental
intervention for conservationists. This is not to say that interventions
need always be cause-effect studies or that evidence-based decisions
cannot be made with compelling preliminary evidence or mensurative data.
We are simply proposing that a statement of the relative strength of
evidence and gaps in the research provides a future direction for
additional research and for implementation.

\textbf{7. Explore the benefits of minimal intervention for
stakeholders.} Resources are limiting, and at times, the
business-as-usual model can provide a guide to intervention for some
environmental management challenges {[}27,28{]}. At the minimum,
exploration of a hope for the best strategy or minimal intervention is
critical because of costs. Business-as-usual models can also provide an
economic mechanism to value ecosystems services {[}29,30{]}, and whilst
this is not without debate, this can expand the breath of stakeholders
and potential investors in a solution for a particular challenge. A best
and worst case scenario analysis is also likely a frequent need for many
environmental challenges because of inertia in the structures that we
use to manage people and resources.

\textbf{8. List the tools applied to this challenge.} In an
environmental management challenge case study, there is typically at
least one primary tool that the researchers used to explore a challenge,
and there are many (Table 1). The tools in basic biology and ecology
relevant to environmental management can be reproducible it at least
conceptually, they can be replicated in another system or applied to
similar challenge - i.e.~citizen science as a means to collect
environmental data {[}31{]} is relevant to many of the challenges we
face including global warming, water quality, and declining
biodiversity.

\textbf{9. Link the primary reproducible tool to the outcome.} The tool
can provide a means to collect data, detect patterns, directly solve an
environmental challenge, or inform policy. If the paper was a direct
test of basic ecology for an environmental challenge, this can be very
straightforward. For instance, ``Odonata (Insecta) as a tool for the
bio-monitoring of environmental quality'' {[}32{]} clearly provides a
means to measure and detect. However, the other proposed roles can
address challenges in a diversity of ways. The identification of or
provision of research evidence is the most ``basic'' role, and it is
also likely the most typical in ecology. Tools that can function in this
capacity include surveys, meta-analyses, systematic reviews, citizen
science data collection, mapping, open-access data, modelling to predict
changes. Tests in the second category directly examine the efficacy of a
management strategy or intervention. Tools that could play this role
include bio-monitoring, mitigation and remediation experiments, and
population demography studies. The inform policy research is typically
more indirect and synthetic. They often take the form of anthropocentric
studies that consider ecological or environmental policy. Any of the
above tools can serve this role, but some tools fit most squarely in
this role include economic incentivization models ({[}33{]}), human
health impact studies {[}34{]}, and human well-being monitoring
associated with environmental interventions {[}35{]}.

\textbf{10. Apply the tool to another challenge or explain how it is
generalizable.} This principle proposes that the primary tool is
reproducible if it can be applied to another challenge or context. It
ties together the concept that reverse-engineered reproducible solutions
are relevant to more than the unpacked, single environmental management
challenge case. This can promote increased in efficiency for tackling
novel environmental challenges as they emerge, and it also supports the
overarching philosophy here for basic science that we cannot continue to
ignore reuse given the global environmental needs for better decision
making.

\section{Implications}\label{implications}

At its core, these principles can distribute the burden of scientific
communication between scientists and stakeholders and embodies a spirit
of dialog between senders and receivers (or between producers and
consumers). Reuse is also not the sole criterion for useful science nor
should it be. Professional advocacy and knowledge mobilization are
increasingly important priorities for universities and science in
general {[}36{]}. Evidence-based decision making is a critical area for
growth and knowledge in many disciplines {[}37--39{]} - not just
environmental management. Increased consumption of scientific evidence
by managers and basic science that is more palatable to a broader
audience written by researchers will result in increased functional use
of scientific literature. Collaboration with stakeholders will
facilitate this process, and open science will be pivotal.

Table 1. Tools for environmental manangement challenges that can be
derived from fundamental research papers.

\section*{References}\label{references}
\addcontentsline{toc}{section}{References}

\hypertarget{refs}{}
\hypertarget{ref-Koontz2018}{}
1. Koontz TM, Thomas CW. Use of science in collaborative environmental
management: Evidence from local watershed partnerships in the Puget
Sound. Environmental Science and Policy. Elsevier; 2018;88: 17--23.
doi:\href{https://doi.org/10.1016/j.envsci.2018.06.007}{10.1016/j.envsci.2018.06.007}

\hypertarget{ref-Cash2003}{}
2. Cash D, Clark W, Alcock F, Dickson N, Eckley N, Guston D, et al.
Knowledge systems for sustainable development. Proceedings of the
National Academy of Sciences of the United States of America. 2003;100:
8086--8091.
doi:\href{https://doi.org/10.1063/1.4795010}{10.1063/1.4795010}

\hypertarget{ref-Iacona2018}{}
3. Iacona GD, Sutherland WJ, Mappin B, Adams VM, Armsworth PR, Coleshaw
T, et al. Standardized reporting of the costs of management
interventions for biodiversity conservation. Conservation Biology.
2018;32: 979--988.
doi:\href{https://doi.org/10.1111/cobi.13195}{10.1111/cobi.13195}

\hypertarget{ref-Naidoo2006}{}
4. Naidoo R, Balmford A, Ferraro PJ, Polasky S, Ricketts TH, Rouget M.
Integrating economic costs into conservation planning. Trends in Ecology
and Evolution. 2006;21: 681--687.
doi:\href{https://doi.org/10.1016/j.tree.2006.10.003}{10.1016/j.tree.2006.10.003}

\hypertarget{ref-Lortie2018}{}
5. Lortie CJ, Filazzola A, Kelsey R, Hart AK, Butterfield HS. Better
late than never: a synthesis of strategic land retirement and
restoration in California. Ecosphere. 2018;9: e02367.
doi:\href{https://doi.org/10.1002/ecs2.2367}{10.1002/ecs2.2367}

\hypertarget{ref-Baron2010}{}
6. Baron N. Escape from the ivory tower: A guide to making your science
matter. Washington, DC: Island Press; 2010.

\hypertarget{ref-Lortie2014}{}
7. Lortie CJ. Formalized synthesis opportunities for ecology: Systematic
reviews and meta-analyses. Oikos. 2014;123: 897--902.
doi:\href{https://doi.org/10.1111/j.1600-0706.2013.00970.x}{10.1111/j.1600-0706.2013.00970.x}

\hypertarget{ref-Lewinsohn2014}{}
8. Lewinsohn TM, Attayde JL, Fonseca CR, Ganade G, Jorge LR, Kollmann J,
et al. Ecological literacy and beyond: Problem-based learning for future
professionals. Ambio. 2014;44: 154--162.
doi:\href{https://doi.org/10.1007/s13280-014-0539-2}{10.1007/s13280-014-0539-2}

\hypertarget{ref-Hao2018}{}
9. Hao J. Reconsidering cause inside the clause in scientific
discourse-from a discourse semantic perspective in systemic functional
linguistics. Text and Talk. 2018;38: 525--550.
doi:\href{https://doi.org/10.1515/text-2018-0013}{10.1515/text-2018-0013}

\hypertarget{ref-Gerstner2017}{}
10. Gerstner K, Moreno-Mateos D, Gurevitch J, Beckmann M, Kamback S,
Jones HP, et al. Will your paper be used in a meta-analysis? Make the
reach of your research broader and longer lasting. Methods in Ecology
and Evolution. 2017;8: 777--784.
doi:\href{https://doi.org/10.1111/2041-210X.12758}{10.1111/2041-210X.12758}

\hypertarget{ref-Baker2016}{}
11. Baker M. Is there a reproducibility crisis? Nature. 2016;533:
452--454.

\hypertarget{ref-Kelly2006}{}
12. Kelly C. Replicating empirical research in behavioral ecology: How
and why it should be done but rarely ever is. The Quarterly Review of
Biology. 2006;81: 221--236.

\hypertarget{ref-Bourne2006}{}
13. Bourne PE, Chalupa LM. Ten simple rules for getting grants. PLoS
Computational Biology. 2006;2: 59--60.
doi:\href{https://doi.org/10.1371/journal.pcbi.0020012}{10.1371/journal.pcbi.0020012}

\hypertarget{ref-Dickman2014}{}
14. Dickman AJ, Hazzah L, Carbone C, Durant SM. Carnivores, culture and
'contagious conflict': Multiple factors influence perceived problems
with carnivores in Tanzania's Ruaha landscape. Biological Conservation.
Elsevier Ltd; 2014;178: 19--27.
doi:\href{https://doi.org/10.1016/j.biocon.2014.07.011}{10.1016/j.biocon.2014.07.011}

\hypertarget{ref-Sandel2015}{}
15. Sandel B. Towards a taxonomy of spatial scale-dependence. Ecography.
2015;38: 358--369.
doi:\href{https://doi.org/10.1111/ecog.01034}{10.1111/ecog.01034}

\hypertarget{ref-Jones2009}{}
16. Jones PG, Thornton PK. Croppers to livestock keepers: livelihood
transitions to 2050 in Africa due to climate change. Environmental
Science and Policy. 2009;12: 427--437.
doi:\href{https://doi.org/10.1016/j.envsci.2008.08.006}{10.1016/j.envsci.2008.08.006}

\hypertarget{ref-Lindsey2009}{}
17. Lindsey PA, Romañach SS, Davies-Mostert HT. The importance of
conservancies for enhancing the value of game ranch land for large
mammal conservation in southern Africa. Journal of Zoology. 2009;277:
99--105.
doi:\href{https://doi.org/10.1111/j.1469-7998.2008.00529.x}{10.1111/j.1469-7998.2008.00529.x}

\hypertarget{ref-Towns2009}{}
18. Towns L, Derocher AE, Stirling I, Lunn NJ, Hedman D. Spatial and
temporal patterns of problem polar bears in Churchill, Manitoba. Polar
Biology. 2009;32: 1529--1537.
doi:\href{https://doi.org/10.1007/s00300-009-0653-y}{10.1007/s00300-009-0653-y}

\hypertarget{ref-Bagchi2006}{}
19. Bagchi S, Mishra C. Living with large carnivores: Predation on
livestock by the snow leopard (Uncia uncia). Journal of Zoology.
2006;268: 217--224.
doi:\href{https://doi.org/10.1111/j.1469-7998.2005.00030.x}{10.1111/j.1469-7998.2005.00030.x}

\hypertarget{ref-Johnson2006}{}
20. Johnson A, Vongkhamheng C, Hedemark M, Saithongdam T. Effects of
human-carnivore conflict on tiger (Panthera tigris) and prey populations
in Lao PDR. Animal Conservation. 2006;9: 421--430.
doi:\href{https://doi.org/10.1111/j.1469-1795.2006.00049.x}{10.1111/j.1469-1795.2006.00049.x}

\hypertarget{ref-Dickman2010}{}
21. Dickman AJ. Complexities of conflict: The importance of considering
social factors for effectively resolving human-wildlife conflict. Animal
Conservation. 2010;13: 458--466.
doi:\href{https://doi.org/10.1111/j.1469-1795.2010.00368.x}{10.1111/j.1469-1795.2010.00368.x}

\hypertarget{ref-Davidson2013}{}
22. Davidson MD. On the relation between ecosystem services, intrinsic
value, existence value and economic valuation. Ecological Economics.
Elsevier B.V. 2013;95: 171--177.
doi:\href{https://doi.org/10.1016/j.ecolecon.2013.09.002}{10.1016/j.ecolecon.2013.09.002}

\hypertarget{ref-Reed2008}{}
23. Reed MS. Stakeholder participation for environmental management: A
literature review. Biological Conservation. 2008;141: 2417--2431.
doi:\href{https://doi.org/10.1016/j.biocon.2008.07.014}{10.1016/j.biocon.2008.07.014}

\hypertarget{ref-Colvin2016}{}
24. Colvin RM, Witt GB, Lacey J. Approaches to identifying stakeholders
in environmental management: Insights from practitioners to go beyond
the 'usual suspects'. Land Use Policy. Elsevier Ltd; 2016;52: 266--276.
doi:\href{https://doi.org/10.1016/j.landusepol.2015.12.032}{10.1016/j.landusepol.2015.12.032}

\hypertarget{ref-Ruxton2017}{}
25. Ruxton G, Colgrave N. Experimental design for the life sciences.
Second. Oxford, UK: Oxford University Press; 2017.

\hypertarget{ref-Kotiaho2002}{}
26. Kotiaho J, Tomkins J. Meta-analysis, can it ever fail? Oikos.
2002;96: 551--553.

\hypertarget{ref-Ferguson2015}{}
27. Ferguson P. The green economy agenda: Business as usual or
transformational discourse? Environmental Politics. 2015;24: 17--37.
doi:\href{https://doi.org/10.1080/09644016.2014.919748}{10.1080/09644016.2014.919748}

\hypertarget{ref-Mosnier2017}{}
28. Mosnier C, Duclos A, Agabriel J, Gac A. What prospective scenarios
for 2035 will be compatible with reduced impact of French beef and dairy
farm on climate change? Agricultural Systems. Elsevier; 2017;157:
193--201.
doi:\href{https://doi.org/10.1016/j.agsy.2017.07.006}{10.1016/j.agsy.2017.07.006}

\hypertarget{ref-Fu2018}{}
29. Fu Q, Hou Y, Wang B, Bi X, Li B, Zhang X. Scenario analysis of
ecosystem service changes and interactions in a mountain-oasis-desert
system: a case study in Altay Prefecture, China. Scientific Reports.
2018;8: 1--13.
doi:\href{https://doi.org/10.1038/s41598-018-31043-y}{10.1038/s41598-018-31043-y}

\hypertarget{ref-Karttunen2018}{}
30. Karttunen K, Ahtikoski A, Kujala S, Törmä H, Kinnunen J, Salminen H,
et al. Regional socio-economic impacts of intensive forest management, a
CGE approach. Biomass and Bioenergy. Elsevier Ltd; 2018;118: 8--15.
doi:\href{https://doi.org/10.1016/j.biombioe.2018.07.024}{10.1016/j.biombioe.2018.07.024}

\hypertarget{ref-McKinley2017}{}
31. McKinley DC, Miller-Rushing AJ, Ballard HL, Bonney R, Brown H,
Cook-Patton SC, et al. Citizen science can improve conservation science,
natural resource management, and environmental protection. Biological
Conservation. Elsevier Ltd; 2017;208: 15--28.
doi:\href{https://doi.org/10.1016/j.biocon.2016.05.015}{10.1016/j.biocon.2016.05.015}

\hypertarget{ref-Miguel2017}{}
32. Miguel TB, Oliveira-Junior JMB, Ligeiro R, Juen L. Odonata (Insecta)
as a tool for the biomonitoring of environmental quality. Ecological
Indicators. Elsevier; 2017;81: 555--566.
doi:\href{https://doi.org/10.1016/j.ecolind.2017.06.010}{10.1016/j.ecolind.2017.06.010}

\hypertarget{ref-Tilman2018}{}
33. Tilman AR, Levin S, Watson JR. Revenue-sharing clubs provide
economic insurance and incentives for sustainability in common-pool
resource systems. Journal of Theoretical Biology. Elsevier Ltd;
2018;454: 205--214.
doi:\href{https://doi.org/10.1016/j.jtbi.2018.06.003}{10.1016/j.jtbi.2018.06.003}

\hypertarget{ref-Chiabai2018}{}
34. Chiabai A, Quiroga S, Martinez-Juarez P, Higgins S, Taylor T. The
nexus between climate change, ecosystem services and human health:
Towards a conceptual framework. Science of the Total Environment. The
Authors; 2018;635: 1191--1204.
doi:\href{https://doi.org/10.1016/j.scitotenv.2018.03.323}{10.1016/j.scitotenv.2018.03.323}

\hypertarget{ref-McKinnon2015}{}
35. McKinnon M. Sustainability: Map the evidence. Nature. 2015;528:
185--187.
doi:\href{https://doi.org/10.1017/CBO9781107415324.004}{10.1017/CBO9781107415324.004}

\hypertarget{ref-Pace2010}{}
36. Pace ML, Hampton SE, Limburg KE, Bennett EM, Cook EM, Davis AE, et
al. Communicating with the public: Opportunities and rewards for
individual ecologists. Frontiers in Ecology and the Environment. 2010;8:
292--298. doi:\href{https://doi.org/10.1890/090168}{10.1890/090168}

\hypertarget{ref-Roy-Byrne2010}{}
37. Roy-Byrne P, Craske MG, Sullivan G, Rose RD, Edlund MJ, Lang AJ, et
al. Delivery of evidence-based treatment for multiple anxiety disorders
in primary care: A randomized controlled trial. JAMA - Journal of the
American Medical Association. 2010;303: 1921--1928.
doi:\href{https://doi.org/10.1001/jama.2010.608}{10.1001/jama.2010.608}

\hypertarget{ref-Tranfield2003}{}
38. Tranfield D, Denyer D, Smart P. Towards a Methodology for Developing
Evidence-Informed Management Knowledge by Means of Systematic Review.
2003;14: 207--222.

\hypertarget{ref-Aarons2011}{}
39. Aarons GA, Hurlburt M, Horwitz SMC. Advancing a conceptual model of
evidence-based practice implementation in public service sectors.
Administration and Policy in Mental Health and Mental Health Services
Research. 2011;38: 4--23.
doi:\href{https://doi.org/10.1007/s10488-010-0327-7}{10.1007/s10488-010-0327-7}

\nolinenumbers


\end{document}

